\documentclass[11pt,a4paper]{article}
\usepackage[T1]{fontenc}
\usepackage[utf8]{inputenc}
\usepackage{babel}
\babelprovide[import,main]{italian}
\usepackage{mathpazo}
\usepackage[scaled=0.92]{helvet}
\renewcommand{\ttdefault}{lmtt}
\usepackage[margin=2.1cm,top=1.8cm,bottom=1.8cm,headsep=0.5cm]{geometry}
\usepackage{enumitem}
\usepackage{xcolor}
\usepackage{titlesec}
\usepackage{fancyhdr}
\usepackage[colorlinks=true,linkcolor=black,urlcolor=blue!50!black]{hyperref}

\setlength{\parindent}{0pt}
\setlength{\parskip}{0.3em}
\setlist{itemsep=0.05em,topsep=0.1em,parsep=0pt}
\titleformat{\section}{\large\bfseries\sffamily}{\thesection}{0.6em}{}
\titlespacing*{\section}{0pt}{0.8em}{0.15em}

\pagestyle{fancy}
\fancyhf{}
\fancyhead[L]{\small\sffamily Lung segmentation hands-on: per i docenti}
\fancyhead[R]{\small\sffamily\thepage}
\renewcommand{\headrulewidth}{0.3pt}

\newcommand{\code}[1]{\texttt{\small #1}}
\newcommand{\btn}[1]{\textsf{\textbf{#1}}}
\newcommand{\goto}{\,$\rightarrow$\,}
\newcommand{\infobox}[2]{%
  \par\vspace{0.3em}\noindent\fcolorbox{black!35}{black!4}{\parbox{\dimexpr\linewidth-2\fboxsep-2\fboxrule}{%
  \textbf{\sffamily #1} #2}}\par\vspace{0.1em}}

\begin{document}

{\sffamily\bfseries\LARGE Il dataset del corso su Kaggle}\par\vspace{0.2em}
{\sffamily\color{black!60} Passi essenziali per caricare \code{luna16-lung-slices} e far girare i notebook}\par
\vspace{0.3em}\noindent\rule{\linewidth}{0.6pt}

\infobox{Serve.}{Un account Kaggle con il \textbf{telefono verificato} (Settings\goto Phone verification):
senza, nei notebook non si attivano né Internet né la GPU. E il file \code{luna16-lung-slices.zip} (396 MB),
\textbf{senza scompattarlo}.}

\section{Creare il dataset}
\begin{enumerate}
\item Su kaggle.com: \btn{Create}\goto\btn{New Dataset}.
\item Trascina lo zip e aspetta che Kaggle finisca di scompattarlo (alcuni minuti).
\item Titolo: \code{luna16-lung-slices}. L'indirizzo mostrato sotto deve finire con \code{/luna16-lung-slices}:
      è il nome della cartella che i notebook cercano.
\item Visibilità \btn{Private} per ora, poi \btn{Create}.
\item Nella pagina del dataset imposta la licenza \textbf{CC BY 4.0}.
\end{enumerate}

\section{Controllare i file}
Nella scheda \btn{Data} devono esserci le cartelle \code{images}, \code{labels}, \code{nodules},
\code{volumes\_lowres}, \code{demo} e i file \code{metadata.csv}, \code{README.md}, \code{LICENSE.txt}.
Dentro le cartelle i file devono finire in \code{.nii.gz}, non in \code{.nii}.

\section{Provare con un notebook}
\begin{enumerate}
\item \btn{Create}\goto\btn{New Notebook}, poi \btn{File}\goto\btn{Import Notebook}: carica \code{00\_environment\_check.ipynb}.
\item Pannello a destra: \btn{Add Input}, cerca \code{luna16-lung-slices}, clicca \btn{+}.
\item \btn{Session options}: Internet \btn{On}, Accelerator \btn{None}.
\item In una cella nuova controlla il percorso (riquadro sotto).
\item Esegui tutto il notebook: deve stampare \code{Libraries: OK} e \code{Dataset: OK} e mostrare una slice CT.
\item Ripeti con \code{solutions/03\_project\_lung\_segmentation.ipynb} e Accelerator \btn{GPU T4 x2}:
      annota il tempo di un'epoca, il Dice di validazione e quello di test.
\end{enumerate}

\infobox{Controllo del percorso.}{Esegui
\par\vspace{0.2em}\hspace*{1em}\code{!find /kaggle/input -maxdepth 4 -name metadata.csv}
\par\vspace{0.2em}Deve stampare
\par\vspace{0.1em}\hspace*{1em}\code{/kaggle/input/luna16-lung-slices/metadata.csv}
\par\hspace*{1em}\code{/kaggle/input/luna16-lung-slices/nodules/metadata.csv}
\par\vspace{0.2em}Se il percorso è diverso, va cambiata la riga \code{DATA\_DIR} in ogni notebook.}

\section{Darlo agli studenti}
\begin{enumerate}
\item \textbf{Pubblico} (consigliato): pagina del dataset\goto\btn{Settings}\goto visibilità \btn{Public}.
\item Oppure \textbf{privato}: \btn{Share} e aggiungi lo username Kaggle di ogni studente.
\item Comunica il nome completo, \code{\textless username\textgreater/luna16-lung-slices}: serve al passo B2 di \code{SETUP.md}.
\end{enumerate}

\section{Se cambiano i dati}
Pagina del dataset\goto\btn{New Version}\goto carica il nuovo zip \textbf{sostituendo} i file.
Nei notebook già aperti: pannello \btn{Input}\goto\btn{Check for updates}. Nome e percorso non cambiano,
quindi i notebook restano uguali.

\infobox{Valori attesi.}{Run di riferimento su CPU: notebook 03, Dice di validazione 0{,}989 e di test 0{,}978;
notebook 04, baseline 0{,}833 e test 0{,}852. Se su GPU escono valori diversi, vanno aggiornati i testi dei notebook.}

\end{document}
